\section{同一技巧在不同 SLO 下方向相反}
\label{sec:slo}

前面三节换过两次切法：按成本链（四类）、按实现平面（三平面）。
本节换第三次：按\textbf{SLO 的类型}。

之所以还要再切一次，是因为它揭示的东西与前两次不同。四类与三平面
回答「有哪些手段」，本节回答\textbf{「同一个手段什么时候是反的」}。

\subsection{一个旋钮，两个 SLO}

最干净的例子是 chunked prefill 的 token 预算。Sarathi-Serve 的原文
把冲突写得很直接：token 预算由\textbf{两个相互竞争的因素}决定——
TBT 的 SLO 要求，与 chunked prefill 自身的开销。\textbf{从压低 TBT
的角度看，预算越小越好}，因为一迭代里 prefill token 越少延迟越低；
但预算太小会导致 prefill 被过度切分，带来开销\texttt{[c-slo]}。

\noindent 这个开销不是含糊的「固定成本」，有确切机制：一个 prompt 被切成
$N$ 个 chunk 时，\textbf{第一个 chunk 的 KV cache 要被加载 $N-1$ 次，
第二个 $N-2$ 次}，以此类推——计算量不变，但 HBM 的重复读取增加
\texttt{[c-slo]}。

\noindent 所以同一个旋钮的两端是这样的：

\begin{center}
\small
\begin{tabular}{@{}p{2.3cm}p{2.2cm}p{2.2cm}@{}}
\toprule
token 预算 & TBT & 吞吐 \\
\midrule
调小 & 变好 & 变差（chunk 过碎，HBM 重复读） \\
调大 & 变差 & 变好 \\
\bottomrule
\end{tabular}
\end{center}

\noindent 这不是「参数有权衡」这种泛泛之谈。它的含义更硬：
\textbf{不存在一个「最优 token 预算」，只存在「对某个 TBT SLO 最优的
预算」}。同一个数在另一个 SLO 下就是错的。

\noindent 同一件事还有一句更直白的表述：优先 prefill 会让已经在跑的
decode 停等，故延迟变高\texttt{[c-slo]}。反过来说，优先 decode
就压住 prefill、抬高 TTFT。\textbf{调度的两端各自伤害一个 SLO}，
没有中间地带能同时最优。

\subsection{应用层：两个 SLO 的相对松紧才是变量}

上面说的是机制。到了部署参数层，问题变成：给定一个应用，
两个 SLO 该怎么配？

DistServe 的三组配置给出了很说明问题的分布\texttt{[c-slo]}：

\begin{table}[t]
\centering
\caption{三组应用的 SLO 配置。TTFT 跨约 120 倍而 TPOT 几乎不变，说明决定配置的是两个 SLO 的相对松紧。}
\label{tab:slo}
\small
\begin{tabular}{@{}lrrl@{}}
\toprule
应用 & TTFT & TPOT & 特征 \\
\midrule
Chatbot OPT-175B & 4.0\,s & 0.2\,s & TTFT 宽松、TPOT 严 \\
Code Completion OPT-66B & 0.125\,s & 0.2\,s & TTFT 极严 \\
Summarization OPT-66B & 15\,s & 0.15\,s & TTFT 极宽松、TPOT 最严 \\
\bottomrule
\end{tabular}
\end{table}

\noindent 看这两个数字的跨度：TTFT 从 $0.125$\,s 到 $15$\,s，
跨约 $120$ 倍；而 TPOT 几乎不变，都在 $0.1$ 到 $0.2$\,s 之间。

\noindent 这个对比说明一件事：\textbf{决定配置的不是「SLO 严不严」，
而是「两个 SLO 的相对松紧」}。代码补全与摘要在「总体严格程度」上
难以比较，但它们的\textbf{偏好方向完全相反}——一个要第一个 token
立刻出来，一个可以等 15 秒但要后续 token 稳定。

\noindent 原文对摘要的判断印证了这一点：摘要任务强调低 TPOT
以便更快生成整段摘要\texttt{[c-slo]}，而它同时容忍 15 秒的 TTFT。

\noindent 这解释了一个实践中反复出现的困惑：为什么同一个分离方案，
在 A 团队那里是 3 倍提升，在 B 团队那里几乎没变。\textbf{不是配置错，
是 SLO 的形状不同。}

\subsection{收益本身随应用变}

DistServe 在同一套系统、同一批硬件上报告了两组差异明显的数字
\texttt{[c-slo]}：

\begin{itemize}[leftmargin=1.3em,itemsep=1.5pt,topsep=2pt]
  \item 代码补全（OPT-66B）：支撑 $3.2\times$ 更高的请求率，
        同时 SLO 更严 $1.5\times$。
  \item 聊天（OPT-13B/66B/175B）：SLO 更严 $1.4\times$ 至 $1.8\times$。
\end{itemize}

\noindent 两者都算「显著提升」，但结构不同：前者主要体现在
\textbf{能扛多少请求}，后者主要体现在\textbf{能扛多严的延迟}。

\noindent 这带来一个方法论上的要求：\textbf{报告分离收益时必须同时
报 SLO 形状}，否则数字不可比。一个「3.2 倍」若不说清是在什么
TTFT/TPOT 组合下达成的，对读者没有决策价值。

\subsection{由此得到的判据}

把本节与前面几节合起来，可以给出一条比第 \ref{sec:boundary} 节
更细的判据：\textbf{先定 SLO 的形状，再定技巧的方向}。

\noindent 具体地，问三个问题：

\begin{enumerate}[leftmargin=1.4em,itemsep=2pt,topsep=3pt]
  \item \textbf{哪个 SLO 更紧？} TTFT 紧（代码补全、搜索）意味着
        优先 prefill、增大 token 预算、倾向不分离或 P 侧多配；
        TPOT 紧（摘要、长文生成）意味着优先 decode、减小预算、
        倾向分离且 D 侧多配。
  \item \textbf{两个 SLO 的比值是多少？} 单独看任何一个都不够。
        摘要的 TTFT 是代码补全的 $120$ 倍，这个比值比绝对值更能
        决定配置。
  \item \textbf{目标达成率是多少？} goodput 的定义是「在 SLO 达成率
        目标（如 $90\%$）下每 GPU 能服务的最大请求率」
        \texttt{[c-slo]}。达成率从 $90\%$ 提到 $99\%$ 会显著改变
        可行配置——尾部延迟的代价不是线性的。
\end{enumerate}

\noindent 这三问的次序不能换。先问「用什么技巧」会得到一堆
互相冲突的建议——因为每个技巧的推荐方向都依赖于 SLO 的形状，
而形状还没定。

\subsection{这一节新增的判断力}

三平面保证不漏，四类给出取舍，而 SLO 是第一个能让前面结论
\textbf{发生方向翻转}的维度——这是它存在的全部理由。
六层的完整关系与使用顺序见第 \ref{sec:onepage} 节的
图 \ref{fig:decision} 与表 \ref{tab:layers}。

